% [原文 p.19]
若 $h'$ 是 $V$ 上的另一度量，则全射同态 $G_n(V_h)\times V\to E$（其中 $E$ 为分类丛）诱导恒等映射 $G_n(V_h)\to G_n(V_{h'})$。这个恒等映射连续，因此 $G_n(V)$ 上的拓扑不依赖于度量的选取。

现在考虑由
\[
(z_1,\dots,z_m)\longmapsto(z_1,\dots,z_{m-1})
\]
给出的自然投影
\[
\C^m\longrightarrow\C^{m-1}.
\]
这个投影诱导连续映射
\[
\iota_{m-1}\colon G_n(\C^{m-1})\longrightarrow G_n(\C^m).
\]
记 $G_n(\C^m)$ 上的分类丛为 $E_{(m)}$，则容易看出
\[
\iota_{m-1}^{*}(E_{(m)})\cong E_{(m-1)}
\]
成立。

\begin{theorem*}{1.4.15}
  对 $f\colon X\to G_n(\C^m)$，由 $f\mapsto f^{*}(E_{(m)})$ 诱导的映射
  \[
  \varinjlim_{m}\,[X,G_n(\C^m)]\longrightarrow\Vect_n(X)
  \]
  对所有紧 Hausdorff 空间 $X$ 都是同构。
\end{theorem*}

% [Tutor Note: 原书 $G_n(V)$ 是 $V$ 的余维 $n$ 子空间全体，不是 $n$ 维子空间全体（定义见原书 p.18）。因此当 $\dim E=n$ 时，$G_n(\C^m)$ 的参数是 $m-n$ 维子空间，也可看成 $n$ 维商空间。若按「$n$ 维子空间」理解定理 1.4.15，维数会对不上。]
\begin{proof}
  构造逆映射。设 $E$ 为 $X$ 上的向量丛。由推论 1.4.13，存在满射同态 $\varphi\colon X\times\C^m\to E$。设 $f\colon X\to G_n(\C^m)$ 为由 $\varphi$ 诱导的映射。若对充分大的 $m'$ 能证明 $f$ 在 $G_n(\C^{m'})$ 中的同伦类不依赖于 $\varphi$ 的选取，则把 $E$ 对应到 $f$ 的同伦类，便构造出逆映射
  \[
  \Vect_n(X)\longrightarrow\varinjlim_{m}\,[X,G_n(\C^m)].
  \]

  设 $\varphi_i\colon X\times\C^{m_i}\to E$（$i=0,1$）为满射同态，$g_i\colon X\to G_n(\C^{m_i})$ 为 $\varphi_i$ 诱导的映射。定义 $\psi_t\colon X\times\C^{m_0}\times\C^{m_1}\to E$ 为
  \[
  \psi_t(x,v_0,v_1)=(1-t)\varphi_0(x,v_0)+t\varphi_1(x,v_1).
  \]
  这是满射同态。设 $f_t\colon X\to G_n(\C^{m_0}\oplus\C^{m_1})$ 为由 $\psi_t$ 诱导的映射。通过
  \[
  (z_1,\dots,z_{m_0})\oplus(u_1,\dots,u_{m_1})\longmapsto(z_1,\dots,z_{m_0},u_1,\dots,u_{m_1})
  \]
% [原文 p.20]
  把 $\C^{m_0}\oplus\C^{m_1}$ 与 $\C^{m_0+m_1}$ 等同起来，则
  \[
  f_0=j_0\circ g_0,\qquad f_1=T\circ j_1\circ g_1
  \]
  成立。其中 $j_i\colon G_n(\C^{m_i})\to G_n(\C^{m_0+m_1})$ 是自然包含映射，
  \[
  T\colon G_n(\C^{m_0+m_1})\longrightarrow G_n(\C^{m_0+m_1})
  \]
  是由 $\C^{m_0+m_1}$ 的坐标置换诱导的映射，同伦于恒等映射。因此 $j_1\circ g_1$ 同伦于 $f_1$，并且如所要求的那样也同伦于 $j_0\circ g_0$。
\end{proof}

\begin{remark}{}
  可以用如下方式把向量丛解释为模。以 $C(X)$ 表示 $X$ 上复值连续函数的环。设 $E$ 是 $X$ 上的向量丛，则 $\Gamma(E)$ 在逐点乘法
  \[
  fs(x)=f(x)s(x)\qquad(f\in C(X),\ s\in\Gamma(E))
  \]
  下成为 $C(X)$ 模。更进一步，同态 $\varphi\colon E\to F$ 给出 $C(X)$ 模同态
  \[
  \Gamma\varphi\colon\Gamma(E)\longrightarrow\Gamma(F).
  \]
  这样，$\Gamma$ 是从 $X$ 上向量丛的范畴 $\cV$ 到 $C(X)$ 模的范畴 $\cM$ 的函子。设 $E$ 是 $n$ 维平凡向量丛，则 $\Gamma(E)$ 是秩 $n$ 的自由模。设 $F$ 也是平凡向量丛，则
  \[
  \Gamma\colon\mathrm{HOM}(E,F)\longrightarrow\Hom_{C(X)}(\Gamma(E),\Gamma(F))
  \]
  是双射。事实上，若选取同构 $E\cong X\times V$，$F\cong X\times W$，则
  \begin{align*}
  \mathrm{HOM}(E,F)&\cong\Hom_{\C}(V,W)^{X}\cong C(X)\otimes\Hom_{\C}(V,W)\\
  &\cong\Hom_{C(X)}(\Gamma(E),\Gamma(F))
  \end{align*}
  成立。因此，$\Gamma$ 诱导出平凡向量丛的范畴 $\cT$ 与有限秩自由 $C(X)$ 模的范畴 $\cF$ 之间的等价。

% [Tutor Note: 这句话只断言 $\Gamma$ 在平凡丛与自由模的子范畴之间是等价；要把等价扩张到全部向量丛，还得靠下一段取投影算子的像所成的子范畴，并配合系 1.4.14。]
% [原文 p.21]
  $\mathrm{Proj}(\cT)$ 取为 $\cV$ 的如下子范畴：其对象是 $\cT$ 中某个投影算子的像。同样地定义 $\mathrm{Proj}(\cF)\subset\cM$。此时容易看出，$\Gamma$ 诱导范畴等价
  \[
  \mathrm{Proj}(\cT)\longrightarrow\mathrm{Proj}(\cF).
  \]
  然而，由推论 1.4.14，$\mathrm{Proj}(\cT)=\cV$。按定义，$\mathrm{Proj}(\cF)$ 是有限生成投射 $C(X)$ 模的范畴。因此，下述命题得证。
\end{remark}

\begin{proposition*}{}
  $\Gamma$ 诱导出 $X$ 上向量丛的范畴与 $C(X)$ 上有限生成投射模的范畴之间的等价。
\end{proposition*}

\section{附加结构}
\takeaway{在向量丛上附加结构，便得到自对偶丛、正交丛、辛丛、自共轭丛、实丛与四元数丛；它们都是更一般的 $F\to G$ 丛的特例，而 $m$ 丛则提示以同伦等价代替同构。}

在线性代数中，常常在向量空间上追加某种结构来考虑问题；对向量丛也可以这样做。例如前面已经讨论过 Hermite 度量。下一个简单的例子，是考虑非退化双线性形式。设 $V$ 为向量丛，$V$ 上的非退化双线性形式是指 $\mathrm{HOM}(V\otimes V,1)$ 中的一个元素 $T$，它对每个 $x\in X$ 都诱导出 $\Hom(V_x\otimes V_x,\C)$ 中的非退化元素。换言之，$T$ 也可以看作 $\mathrm{ISO}(V,V^{*})$ 的元素。把向量丛 $V$ 与这个同构映射 $T$ 合在一起，称为自对偶丛（self-dual bundle）。

% [Tutor Note: 这里的 $1$ 是平凡线丛，不是数字；$\mathrm{HOM}(V\otimes V,1)$ 是 $V\otimes V$ 的对偶丛的截面空间，也就是 $V$ 上双线性形式组成的向量空间。前面各节只用 $\Hom$ 记向量空间之间的同态，此处大写的 $\mathrm{HOM}$ 专指丛同态的截面集合。]
当 $T$ 对称，即对每个 $x\in X$，$T_x$ 都对称时，称 $(V,T)$ 为正交丛（orthogonal bundle）。当 $T$ 斜对称，即对每个 $x\in X$，$T_x$ 都斜对称时，称 $(V,T)$ 为辛丛（symplectic bundle）。

作为另一个例子，也可以考虑 $T\in\mathrm{ISO}(V,\bar V)$ 的对 $(V,T)$。其中 $\bar V$ 表示 $V$ 的复共轭丛（complex conjugate bundle），即对 $V$ 施以「复共轭函子」所得的丛。这样的 $(V,T)$ 可以称为自共轭丛（self-conjugate bundle）。同构映射 $T$ 也可以看作反线性同构 $V\to\bar V$。于是可以作 $T^2$。当 $T^2=\mathrm{id}_V$ 时，称 $(V,T)$ 为实向量丛\footnote{译注：关于实向量丛，参见附录 C“K 理论与实结构”。（原日译本 p.22 译注）}。

% [原文 p.22]
事实上，由所有满足 $Tv=v$ 的 $v\in V$ 构成的子空间 $W\subset V$ 具有这样的向量丛结构：每条纤维都是实数上的向量空间；而 $V$ 可以等同于 $W$ 的复化 $W\otimes_{\R}\C$。当 $T^2=-\mathrm{id}_V$ 时，称 $(V,T)$ 为四元数丛。事实上，$V_x$ 上的四元数向量空间结构可由 $j(v)=Tv$ 定义。这里四元数由 $i,j$ 生成，满足 $ij=-ji$，且在 $\R$ 上 $i^2=j^2=-1$。

% [Tutor Note: 自共轭丛的 $T$ 是反线性同构，故 $T^2$ 才是线性的自同态；$T^2=\mathrm{id}_V$ 给出实结构，$T^2=-\mathrm{id}_V$ 给出四元数结构，这两个符号约定在第 2 章讨论 $K$ 理论中的实结构时会反复用到。]
如果 $V$ 带有 Hermite 度量 $h$，则它给出同构映射 $\bar V\to V^{*}$，把自共轭丛对应到自对偶丛。细致考察对称与斜对称两种情形，并证明自共轭丛、正交丛、辛丛分别与自对偶丛、实向量丛、四元数丛在相差一个同伦的意义下本质等价，这作为习题留给读者\footnote{关键在于 $\mathrm{GL}(n,\R)$ 与 $\mathrm{O}(n,\C)$ 具有相同的极大紧子群 $\mathrm{O}(n,\R)$。同样的注意事项对斜对称的情形也适用。}。因此，可以按考察对象选取方便的一侧。例如，把前面各节的结果推广到实向量丛或四元数丛是容易的，而把引理 1.4.3 推广到正交丛或辛丛就不那么简单。另一方面，张量积的性质大概更适合用双线性形式的框架来处理。$(V,T)$ 与 $(W,S)$ 的张量积是 $(V\otimes W,T\otimes S)$，而 $T\otimes S$ 的对称性可以由 $T$ 与 $S$ 的对称性立即推出。特别要注意，两个辛丛的张量积是正交丛。

自共轭丛是更一般概念的特别情形。设 $F$，$G$ 都是向量空间上的连续函子。此时，若 $V$ 是向量丛且 $T\in\mathrm{ISO}(F(V),G(V))$，则把对 $(V,T)$ 记作 $F\to G$ 丛。显然，取 $F(V)=V$，$G(V)=\bar V$ 就得到自共轭丛。另一个重要的例子是：固定 $m$，取 $F$ 与 $G$ 为如下的 $m$ 重函子：
\[
F(V)=G(V)=\underbrace{V\oplus V\oplus\cdots\oplus V}_{m\ \text{个}},
\]
于是 $m\to m$ 丛（或更简单地，$m$ 丛）就是满足 $T\in\Aut(mV)$ 的对 $(V,T)$。若存在 $S\in\Aut(V)$ 使 $T=mS$，则 $m$ 丛 $(V,T)$ 是平凡的。

% [原文 p.23]
一般地，对 $F\to G$ 丛，与引理 1.4.3 类似的结果不成立，即同伦不能推出同构。因此，要想适当地定义等价的概念，必须把同伦也纳入进来。例如，当存在 $X\times I$ 上的 $m$ 丛 $(W,S)$ 使得
\[
(V_i,T_i)\cong(W,S)|_{X\times\{i\}}\qquad(i=0,1)
\]
时，就称 $m$ 丛 $(V_0,T_0)$ 与 $(V_1,T_1)$ 等价。

\begin{remark}{}
  $X$ 上的 $m$ 丛应当看作 $S(X)$ 上的「模 $m$ 的向量丛」。
\end{remark}

% [Tutor Note: 原书 p.23 此处作 $S(X)$，但本节及前面各节都没有定义过 $S(X)$；按上下文，这句的意思应是 $m$ 丛应与 $K$ 群模 $m$ 的对象对应，$S(X)$ 疑为 $K(X)$ 或同类记号的笔误。]

\section{$G$ 空间上的 $G$-向量丛}
\takeaway{$G$ 向量丛同时推广了普通向量丛与有限群的表示。平均化算子给出不变截面与不变计量，$G$ 同伦不变性、$G$ 向量丛的直和分解以及与投射模的对应都由它导出。}

设 $G$ 为拓扑群。此时，$G$ 空间是指拓扑空间 $X$ 连同 $G$ 在 $X$ 上的连续作用；即 $G$ 作用于 $X$，且映射 $G\times X\to X$ 连续。$G$ 空间之间的 $G$ 映射是指与 $G$ 的作用可交换的映射。$G$ 空间 $E$ 称为 $G$ 空间 $X$ 上的 $G$ 向量丛，是指满足以下条件：

\begin{enumerate}[label=(\roman*),leftmargin=2.2em,itemsep=2pt]
\item $E$ 是 $X$ 上的向量丛；
\item 投影 $E\to X$ 是 $G$ 映射；
\item 对每个 $g\in G$，映射 $E_x\to E_{g(x)}$ 是向量空间的同态。
\end{enumerate}

若 $G$ 是平凡群，则当然所有空间都是 $G$ 空间，所有向量丛都是 $G$ 向量丛。另一个极端是 $X$ 为单点：此时对任何 $G$，$X$ 都是 $G$ 空间，而 $X$ 上的 $G$ 向量丛正是 $G$ 的（有限维）表示空间。因此，$G$ 向量丛是同时包含通常的向量丛与 $G$ 模的自然推广。紧空间上向量丛理论的大部分内容，在 $G$ 也紧时可以推广到 $G$ 向量丛。然而，这一推广以紧群表示的基本事实为前提。所以，这里只限于不含分析学问题的有限群。

% [原文 p.24]
两种极端的情形是如下的 $G$ 空间：

\begin{enumerate}[label=(\roman*),leftmargin=2.2em,itemsep=2pt]
\item $g\neq 1\Rightarrow g(x)\neq x$，即 $X$ 是自由 $G$ 空间的情形；
\item 对所有 $x\in X$，$g\in G$，$g(x)=x$，即 $X$ 是平凡 $G$ 空间的情形。
\end{enumerate}

下面考察这两个极端情形下 $G$ 向量丛的结构。

设 $X$ 为自由 $G$ 空间，$X/G$ 为其轨道空间。此时 $\pi\colon X\to X/G$ 是有限覆盖映射。设 $E$ 为 $X$ 上的 $G$ 向量丛。那么 $E$ 必然是自由 $G$ 空间。轨道空间 $E/G$ 带有 $X/G$ 上的自然向量丛结构。事实上，$E/G\to X/G$ 与 $E\to X$ 局部同构，因此只要 $E$ 局部平凡，$E/G$ 也局部平凡。反之，设 $V$ 为 $X/G$ 上的向量丛，则 $\pi^{*}V$ 是 $X$ 上的 $G$ 向量丛。事实上，$\pi^{*}V\subset X\times V$，而 $g(x,v)=(g(x),v)$ 定义了 $G$ 在 $X\times V$ 上的作用。显然，对应 $E\mapsto E/G$ 与 $V\mapsto\pi^{*}V$ 互为逆函子。这样就得到下述命题。

\begin{proposition*}{1.6.1}
  若 $X$ 是自由 $G$ 空间，则 $X$ 上的 $G$ 向量丛与 $X/G$ 上的向量丛通过 $E\mapsto E/G$ 一一对应。
\end{proposition*}

在讨论平凡 $G$ 空间之前，回忆有限群表示论的一个基本事实：存在 $G$ 的不可约表示的有限集合 $V_1,\dots,V_k$，使得 $G$ 的任意表示 $V$ 都能唯一地分解为直和 $\sum_{i=1}^{k}n_iV_i$。现在，对任意两个 $G$ 模（即表示空间）$V$，$W$，可以定义 $G$ 同态映射的向量空间 $\Hom_G(V,W)$。此时
\[
\Hom_G(V_i,V_j)\cong
\begin{cases}
0 & (i\neq j\ \text{的情形}),\\
\C & (i=j\ \text{的情形})
\end{cases}
\]
成立。因此，对任意 $V$，自然映射
% [原文 p.25]
\[
\sum_{i=1}^{k}V_i\otimes\Hom_G(V_i,V)\longrightarrow V
\]
是 $G$ 同构映射。一旦写成这个形式，就可以把该结果推广到平凡 $G$ 空间上的 $G$ 向量丛。事实上，设 $E$ 为平凡 $G$ 空间 $X$ 上的任意 $G$ 向量丛，则可以定义自同态映射 $\mathrm{Av}\in\End(E)$ 为
\[
\mathrm{Av}(e)=\frac{1}{|G|}\sum_{g\in G}g(e)\qquad(e\in E),
\]
其中 $|G|$ 表示 $G$ 的阶（这里用到了如下事实：当 $X$ 是平凡 $G$ 空间时，每个 $g\in G$ 都定义了 $E$ 的一个自同态映射）。容易看出 $\mathrm{Av}$ 是 $E$ 上的投影算子，因此其像是 $E$ 的不变子空间，并且是向量丛。记之为 $E^{G}$，称为 $E$ 的不变子丛（invariant sub-bundle）。这样，若 $E$，$F$ 都是 $G$ 向量丛，则 $\Hom_G(E,F)=(\Hom(E,F))^{G}$ 也是向量丛。

特别地，取 $E$ 为具有自然 $G$ 作用的平凡丛 $\mathbf{V}_i=X\times V_i$，则可以考虑自然的向量丛同态
\[
\sum_{i=1}^{k}\mathbf{V}_i\otimes\Hom_G(\mathbf{V}_i,F)\longrightarrow F.
\]
对 $G$ 模 $F$，已知它是 $G$ 同构映射。换言之，对 $X$ 上的任意 $G$ 向量丛 $F$，它对每个 $x\in X$ 都是 $G$ 同构映射。因此，它是 $G$ 向量丛的同构映射。所以，每个 $G$ 向量丛 $F$ 都同构于形如 $\sum\mathbf{V}_i\otimes E_i$ 的 $G$ 向量丛，其中 $E_i$ 是具有平凡 $G$ 作用的向量丛。此外，$E_i$ 在差一个同构的意义下唯一确定。事实上，
\begin{align*}
\Hom_G(V_i,F)&\cong\sum_{j=1}^{k}\Hom_G(V_i,V_j\otimes E_j)\\
&\cong\sum_{j=1}^{k}\Hom_G(V_i,V_j)\otimes E_j\\
&\cong E_i
\end{align*}
成立。这样，下述命题得证。

% [Tutor Note: 粗体的 $\mathbf{V}_i=X\times V_i$ 是平凡 $G$ 向量丛，斜体的 $V_i$ 是不可约 $G$ 模；上面这个同构把丛的分解化归为表示的分解，粗体与斜体混用时须分清。]
% [原文 p.26]
\begin{proposition*}{1.6.2}
  设 $X$ 为平凡 $G$ 空间，$V_1,\dots,V_k$ 为不可约 $G$ 模的完全代表系，$\mathbf{V}_i=X\times V_i$ 为对应的 $G$ 向量丛。则 $X$ 上的每个 $G$ 向量丛 $F$ 都同构于直和 $\sum_{i=1}^{k}\mathbf{V}_i\otimes E_i$，其中 $E_i$ 是具有平凡 $G$ 作用的向量丛。此外，$E_i$ 在差一个同构的意义下唯一确定，且由 $E_i=\Hom_G(\mathbf{V}_i,F)$ 给出。
\end{proposition*}

现在回到一般的（紧）$G$ 空间 $X$ 的情形，说明 1.4 节的结果如何推广到 $G$ 向量丛。首先注意，若 $E$ 是 $G$ 向量丛，则
\[
(gs)(x)=g(s(g^{-1}(x)))\qquad(s\in\Gamma(E))
\]
使 $G$ 自然作用于 $\Gamma(E)$。截面 $s$ 是不变的，是指对所有 $g\in G$ 都有 $gs=s$。所有不变截面的集合构成 $\Gamma(E)$ 的子空间 $\Gamma(E)^{G}$。如通常那样，平均化算子
\[
\mathrm{Av}=\frac{1}{|G|}\sum_{g\in G}g
\]
定义了在 $\Gamma(E)^{G}$ 上为恒等映射的同态 $\Gamma(E)\to\Gamma(E)^{G}$。

\begin{lemma*}{1.6.3}
  设 $X$ 为紧 $G$ 空间，$Y\subset X$ 为闭子 $G$ 空间（即 $G$ 作用下的不变闭子空间），$E$ 为 $X$ 上的 $G$ 向量丛。此时，任意不变截面 $s\colon Y\to E|_Y$ 都可延拓为 $X$ 上的不变截面。
\end{lemma*}

% [Tutor Note: 平均化算子要除以 $|G|$，所以在特征不整除 $|G|$ 的域上才可逆；本书只在复系数上工作，有限群这一限制正是为了让平均化恒可行。若换成紧群，则应把求和换成 Haar 积分。]
\begin{proof}
  由引理 1.4.1，可把 $s$ 延拓为 $X$ 上的向量丛 $E$ 的截面 $t$。此时 $\mathrm{Av}(t)$ 是 $X$ 上的向量丛 $E$ 的不变截面，而在 $Y$ 上 $s$ 不变，故
  \[
  \mathrm{Av}(t)=\mathrm{Av}(s)=s
  \]
  成立。因此 $\mathrm{Av}(t)$ 即为所求的延拓。
\end{proof}

设 $E$，$F$ 为 $G$ 向量丛，则 $\Hom(E,F)$ 也是 $G$ 向量丛，且
\[
\Gamma(\Hom(E,F))^{G}\cong\mathrm{HOM}_G(E,F)
\]
成立。

% [原文 p.27]
因此，由引理 1.6.3 立即推出引理 1.4.2 与引理 1.4.3 的 $G$ 同变版本。这样，下述引理得证。

\begin{lemma*}{1.6.4}
  设 $Y$ 为紧 $G$ 空间，$X$ 为 $G$ 空间，$f_t\colon Y\to X$（$0\le t\le 1$）为 $G$ 同伦，$E$ 为 $X$ 上的 $G$ 向量丛。此时，$f_0^{*}E$ 与 $f_1^{*}E$ 是同构的 $G$ 向量丛。
\end{lemma*}

$G$ 同伦当然是指 $G$ 映射 $F\colon Y\times I\to X$，其中 $I$ 是带平凡 $G$ 作用的单位区间。$G$ 空间称为 $G$ 可缩的，是指它 $G$ 同伦等价于单点。特别地，$G$ 空间上的锥总是 $G$ 可缩的。平凡 $G$ 向量丛是指：取 $G$ 模 $V$，与积 $X\times V$ 同构的 $G$ 向量丛。按这些定义，引理 1.4.4 到推论 1.4.11 都原样推广到 $G$ 向量丛。只需注意：若 $h$ 是 $E$ 的度量，则 $\mathrm{Av}(h)$ 是不变度量。

要推广引理 1.4.12，注意：若 $V\subset\Gamma(E)$ 丰富，则 $\sum_{g\in G}gV\subset\Gamma(E)$ 既丰富又不变。由此立即得到推论 1.4.14 的适当推广。

推广定理 1.4.15 时，必须考虑 $m\to\infty$ 时 $m\sum_{i=1}^{k}V_i$ 的 $G$ 子空间的 Grassmann 流形，其中 $V_1,\dots,V_k$ 一如既往表示不可约 $G$ 模的完全代表系。这个表述留给读者。

最后，考虑向量丛作为模的解释。设 $A=C(X)$。此时，若 $X$ 是 $G$ 空间，则 $G$ 作为代数自同构的群作用于 $A$。若 $E$ 是 $X$ 上的 $G$ 向量丛，则 $\Gamma(E)$ 是投射 $A$ 模，$G$ 作用于 $\Gamma(E)$，且 $A$ 作用与 $G$ 作用之间的关系如下：
\[
g(as)=g(a)g(s)\qquad(a\in A,\ g\in G,\ s\in\Gamma(E)).
\]
这个关系式可以通过在 $A$ 上引入 $G$ 的「扭群环」$B$ 而获得另一种看法。即，$B$ 的元素是线性组合 $\sum_{g\in G}a_g g$（$a_g\in A$），积由
\[
(ag)(a'g')=(ag(a'))gg'
\]
定义。

% [原文 p.28]
于是 $\Gamma(E)$ 成为 $B$ 模。证明 $X$ 上的 $G$ 向量丛的范畴与有限生成且在 $A$ 上投射的 $B$ 模的范畴等价，作为习题留给读者。

% [Tutor Note: 这个 $B$ 就是交叉积 $A\rtimes G$（扭群环）：乘法里 $g$ 先作用到系数 $a'$ 上再相乘，因此 $B$ 既不是 $A$ 与群环 $\C[G]$ 的张量积，也不是直积；称它为「扭」群环正是指这种系数被 $G$ 扭一下的乘法。]

\chapterreview{%
第 1 章是全书的入口。它一个 $K$ 群也不谈，却把向量丛当作独立的研究对象，系统地建立后面两章反复调用的整套“几何词典”：丛的运算、丛的同态、同伦分类以及与模的对应。Atiyah 的写法是先给出局部定义（向量空间族），再补上整体条件（局部平凡），最后用紧性把局部信息粘成整体结论；这一套路在第 2、3 章一再重演。

全章沿一条单线推进：§1.1–§1.3 建立范畴语言（子丛、商丛、丛同态与直和分解）；§1.4 处理紧空间上的整体理论（同伦不变性、截面延拓与分类定理）；§1.5 讨论附加结构（自对偶、正交、辛、自共轭、实、四元数六种丛以及 $m$ 丛）；§1.6 把整套理论搬到 $G$ 空间上。前半属拓扑，后半属表示论，两者在 $G$ 向量丛的直和分解定理处汇合。

核心结论可归为三条。其一，紧 Hausdorff 空间上的 $n$ 维复向量丛由稳定同伦类分类：$\Vect_n(X)\cong\varinjlim_{m}[X,G_n(\C^m)]$，这是第 2 章定义约化 $K$ 群的出发点。其二，向量丛的范畴与有限生成投射模的范畴等价（Serre–Swan 定理），丛的同态就是模同态。其三，$G$ 向量丛在自由 $G$ 空间与平凡 $G$ 空间两种极端下都有清晰的结构定理，后者给出 $F\cong\sum_i\mathbf{V}_i\otimes E_i$。后两条是同一想法在不同系数环上的展开。

几处约定极易出错。$G_n(V)$ 是余维 $n$ 子空间全体，不是 $n$ 维子空间全体；分类定理中的极限是正向极限，$G_n(\C^m)$ 通过自然包含进入 $G_n(\C^{m+1})$；平均化算子只在有限群（或紧群）上可用，这正是 §1.6 把 $G$ 限制为有限群的理由；§1.5 反复出现的“相差一个同伦”是不可省略的限定词，去掉它，正交丛与实丛并不等价。此外，$m$ 丛的等价关系本身是按同伦定义的，与通常的丛同构是两回事。

本章没有出现一个 $K$ 群，却处处为 $K$ 理论作准备：直和、张量积、对偶都在这里定义；同伦不变性与“丛的范畴等价于模的范畴”这两件事，正是第 2 章定义 $K(X)$、并讨论其周期性的出发点。读者若把本章当作拓扑复习，读到第 2 章会觉得突兀；若把它当作 $K$ 群的几何来源，第 2 章的定义几乎是被逼出来的。}

% CHUNK ch01c DONE
